% Part: turing-machines
% Chapter: machines-computations
% Section: combining-machines

\documentclass[../../../include/open-logic-section]{subfiles}

\begin{document}

\olfileid{tur}{mac}{cmb}
\olsection{Nggabungake Mesin Turing}

\begin{explain}
Tuladha mesin Turing kang wis kita deleng nganti saiki
lumayan prasaja. Nanging satemene, saben masalah kang
bisa dirampungake nganggo basa pamrograman modern uga
bisa dirampungake nganggo mesin Turing. Kanggo mbangun
mesin Turing kang luwih rumit, kita perlu mesthekake
yen mesin-mesin iku bisa digabung. Kanthi mangkono,
kita bisa mbangun mesin kanggo ngrampungake masalah
luwih rumit kanthi mecah prosedure dadi bagean-bagean
luwih prasaja. Yen kita bisa nemokake cara lumrah
kanggo mecah masalah rumit dadi bagean-bageane,
kita bisa ngrampungake masalah iku ing sawetara
tahap: gawe sawetara mesin Turing prasaja banjur
nggabungake dadi siji mesin kang bisa ngrampungake
masalah mau. Bab iki mligine wigati nalika kita
ngadhepi Masalah Mandheg ing bagean sabanjure.

Kepriye carane nggabungake mesin Turing
$M = \tuple{Q, \Sigma, q_0, \delta}$ lan~$M' =
\tuple{Q', \Sigma', q_0', \delta'}$? Konfigurasi pita
sawise $M$~mandheg saiki kita gunakake minangka
konfigurasi input kanggo lumakune mesin~$M'$.
Kanggo ngasilake siji mesin Turing $M \frown M'$
kang nindakake iki, tindakna langkah-langkah iki:
\begin{enumerate}
    \item Wenehana nomer (utawa jeneng) anyar marang kabeh
    kahanan~$Q'$ saka~$M'$, supaya $M$ lan~$M'$
    ora nduwé kahanan kang padha ($Q \cap Q' = \emptyset$).
    \item Himpunan kahanan $M \frown M'$ yaiku $Q \cup Q'$.
    \item Alfabet pitane yaiku $\Sigma \cup \Sigma'$.
    \item Kahanan wiwitane yaiku~$q_0$.
    % OLPL-149: Ing sumber beku, syarat cabang pisanan uga nyakup transisi kang ora ditegesi; syarat iki diwatesi supaya ora tabrakan karo cabang katelu.
    \item Fungsi transisine yaiku fungsi $\delta''$ iki:
    \[\delta''(q,\sigma) =
    \begin{cases}
      \delta(q,\sigma) & \text{yen $q \in Q$ lan transisine ditegesi}\\
      \delta'(q,\sigma) & \text{yen $q \in Q'$}\\
      \tuple{q_0', \sigma, \TMstay} & \text{yen $q \in Q$ lan
      $\delta(q,\sigma)$ ora ditegesi}
    \end{cases}\]
\end{enumerate}
Mesin kang diasilake nggunakake instruksi saka~$M$
nalika ana ing kahanan $q \in Q$, lan instruksi
saka~$M'$ nalika ana ing kahanan~$q \in Q'$.
Nalika ana ing kahanan $q \in Q$ lan lagi maca
simbol~$\sigma$ kang ora nduwé transisi ing~$M$
(yaiku nalika $M$ mesthine mandheg), mesin mau
mlebu kahanan wiwitan~$M'$ (tanpa ngowahi isi
pita utawa panggonan sirah maca-nulis).

Elinga, yen mesin~$M$ ora tertib, kita ora ngerti
ing endi sirah maca-nulis nalika $M$~mandheg.
Mula, ing konfigurasi nalika $M$ mandheg, sirah
durung mesthi maca kothak~$1$. Iki kudu digatekake
nalika nggabungake mesin.
\end{explain}

\begin{ex}
\emph{Nggabungake Mesin:} Kita bakal ngrancang mesin kang,
yen diwiwiti nganggo input rong blok~$\TMstroke$ kang
dawane $n$ lan~$m$, mandheg kanthi siji blok
$2(m+n)$~$\TMstroke$ ing pita. Kanggo mbangun mesin
iki, kita bisa nggabungake rong mesin kang wis kita
kenal: mesin panambah lan mesin pangganda. Ayo diwiwiti
kanthi nggambar diagram kahanan mesin panambah.
\[
\begin{tikzpicture}[->,>=stealth',shorten >=1pt,auto,node distance=2.8cm,
                    semithick]
  \tikzstyle{every state}=[fill=none,draw=black,text=black]

  \node[initial,state] (A)              {$q_0$};
  \node[state]         (B) [right of=A] {$q_1$};
  \node[state]         (C) [right of=B] {$q_2$};

  \path (A) edge node {\TMtrans{\TMblank}{\TMstroke}{\TMstay}} (B)
            edge [loop above] node {\TMtrans{\TMstroke}{\TMstroke}{\TMright}} (B)
        (B) edge [loop above] node {\TMtrans{\TMstroke}{\TMstroke}{\TMright}} (B)
            edge node {\TMtrans{\TMblank}{\TMblank}{\TMleft}} (C)
        (C) edge [loop above] node {\TMtrans{\TMstroke}{\TMblank}{\TMstay}} (C);
\end{tikzpicture}
\]
Tinimbang mandheg ing kahanan~$q_2$, kita kepengin
mesin nerusake pakaryane kanggo nggandakake output.
Elinga, mesin pangganda mbusak guratan kapisan ing
input lan nulis rong guratan ing output kang kapisah.
Ayo nambah instruksi supaya sirah maca-nulis mesthi
maca guratan kapisan saka output mesin panambah.
\[
\begin{tikzpicture}[->,>=stealth',shorten >=1pt,auto,node distance=2.8cm,
                    semithick]
  \tikzstyle{every state}=[fill=none,draw=black,text=black]

  \node[initial,state] (A)              {$q_0$};
  \node[state]         (B) [right of=A] {$q_1$};
  \node[state]         (C) [right of=B] {$q_2$};
  \node[state]         (D) [below left of=C] {$q_3$};
  \node[state]         (E) [below left of=D] {$q_4$};

  \path (A) edge node {\TMtrans{\TMblank}{\TMstroke}{\TMstay}} (B)
            edge [loop above] node {\TMtrans{\TMstroke}{\TMstroke}{\TMright}} (B)
        (B) edge [loop above] node {\TMtrans{\TMstroke}{\TMstroke}{\TMright}} (B)
            edge node {\TMtrans{\TMblank}{\TMblank}{\TMleft}} (C)
        (C) edge node {\TMtrans{\TMstroke}{\TMblank}{\TMleft}} (D)
        (D) edge [loop left] node {\TMtrans{\TMstroke}{\TMstroke}{\TMleft}} (D)
            edge node[left, xshift=-2mm] {\TMtrans{\TMendtape}{\TMendtape}{\TMright}} (E);
\end{tikzpicture}
\]
Saiki gampang kanggo nggandakake input---kita mung
perlu nyambungake mesin pangganda marang kahanan~$q_4$.
Kanggo iki, kahanan-kahanan mesin pangganda kudu diwenehi
jeneng anyar supaya diwiwiti saka~$q_4$, dudu saka~$q_0$
---kanthi mangkono ora ana rong kahanan wiwitan.
Diagram pungkasan kudune kaya ing \olref{fig:combined}.
\begin{figure}
\[
\begin{tikzpicture}[->,>=stealth',shorten >=1pt,auto,node distance=2.8cm,
                    semithick]
  \tikzstyle{every state}=[fill=none,draw=black,text=black]
  \node[initial,state] (A)              {$q_0$};
  \node[state]         (B) [right of=A] {$q_1$};
  \node[state]         (C) [right of=B] {$q_2$};
  \node[state]         (D) [below left of=C] {$q_3$};
  \node[state]         (E) [below left of=D] {$q_4$};
  \node[state]         (2) [right of=E] {$q_5$};
  \node[state]         (3) [right of=2] {$q_6$};
  \node[state]         (4) [below of=3] {$q_7$};
  \node[state]         (5) [left of=4]  {$q_8$};
  \node[state]         (6) [left of=5]  {$q_9$};

  \path (A) edge node {\TMtrans{\TMblank}{\TMstroke}{\TMstay}} (B)
            edge [loop above] node {\TMtrans{\TMstroke}{\TMstroke}{\TMright}} (B)
        (B) edge [loop above] node {\TMtrans{\TMstroke}{\TMstroke}{\TMright}} (B)
            edge node {\TMtrans{\TMblank}{\TMblank}{\TMleft}} (C)
        (C) edge node {\TMtrans{\TMstroke}{\TMblank}{\TMleft}} (D)
        (D) edge [loop left] node {\TMtrans{\TMstroke}{\TMstroke}{\TMleft}} (D)
            edge node[left, xshift=-2mm] {\TMtrans{\TMendtape}{\TMendtape}{\TMright}} (E)
    (E) edge node {\TMtrans{\TMstroke}{\TMblank}{\TMright}} (2)
    (2) edge [loop above] node {\TMtrans{\TMstroke}{\TMstroke}{\TMright}} (2)
      edge node {\TMtrans{\TMblank}{\TMblank}{\TMright}} (3)
    (3) edge [loop above] node {\TMtrans{\TMstroke}{\TMstroke}{\TMright}} (3)
        edge node {\TMtrans{\TMblank}{\TMstroke}{\TMright}} (4)
    (4) edge [loop below] node {\TMtrans{\TMblank}{\TMstroke}{\TMleft}} (4)
        edge node {\TMtrans{\TMstroke}{\TMstroke}{\TMleft}} (5)
    (5) edge [loop below]  node {\TMtrans{\TMstroke}{\TMstroke}{\TMleft}} (5)
        edge              node {\TMtrans{\TMblank}{\TMblank}{\TMleft}} (6)
    (6) edge [loop below] node {\TMtrans{\TMstroke}{\TMstroke}{\TMleft}} (6)
        edge              node {\TMtrans{\TMblank}{\TMblank}{\TMright}} (E);
\end{tikzpicture}
\]
\caption{Nggabungake mesin panambah lan pangganda}
\ollabel{fig:combined}
\end{figure}
\end{ex}

\begin{prop}
    Yen $M$ lan $M'$ padha-padha tertib lan ngitung fungsi $f\colon
    \Nat^k \to \Nat$ lan $f'\colon \Nat \to \Nat$, miturut urutane,
    mula $M \frown M'$ uga tertib lan ngitung~$\comp{f}{f'}$.
\end{prop}

\begin{proof}
    Amarga $M$ tertib, nalika mandheg kanthi
    output~$f(n_1,\dots,n_k) = m$, sirah maca-nulis
    lagi maca kothak~$1$. Yen banjur mlebu kahanan
    wiwitan~$M'$, mesin bakal mandheg kanthi output
    $f'(m)$ lan sirah maca-nulis bali maca kothak~$1$.
    Syarat liyane saka \olref[dis]{defn:disciplined}
    uga katepung.
\end{proof}

\begin{prob}
    Wenehana mesin Turing tertib kang ngitung $f(x) = x+2$
    kanthi njupuk mesin~$M$ saka
    \cref{tur:mac:dis:prob:disc-succ} lan mbangun $M \frown M$.
\end{prob}
\end{document}
